
\section{Discussion}
\label{sec:discussion}

\subsection{Interpretation of RQ2}
\label{subsec:discussion-rq2}

The synthetic RQ2 analysis produced higher
emotional-dependence marker frequencies in the
long-duration group than in the short-duration group,
as shown in Table~\ref{tab:rq2-summary}.

The Mann--Whitney test, Cliff's delta, and the
trend test identified differences in the generated
data. This demonstrates that the analytical pipeline
can compute group comparisons, effect sizes, and
confidence intervals.

However, these patterns originate from synthetic
inputs and cannot establish an association between
AI-companion use and emotional dependence among
real students.

\subsection{Interpretation of RQ3}
\label{subsec:discussion-rq3}

The synthetic RQ3 analysis produced a weak positive
Spearman correlation between interaction duration
and the social-avoidance index.

However, the adjusted OLS duration coefficient was
not statistically significant, and the model explained
less than 1\% of the outcome variation
($R^2=0.008$).

Spearman correlation measures rank-based association,
whereas the OLS coefficient estimates a conditional
linear relationship after accounting for the included
covariate. Their different results may reflect the
distinct statistical quantities being estimated.

These findings demonstrate the importance of
distinguishing statistical significance from
effect magnitude and practical relevance.

\subsection{Relation to Prior Research}
\label{subsec:prior-research}

Previous studies have investigated AI companions
in relation to loneliness and perceived emotional
support~\cite{defreitasAICompanionsReduce2024,
maplesLonelinessSuicideMitigation2024}.

Other research has examined anthropomorphism,
trust, and relationship development in interactions
with conversational technologies~
\cite{seymourExploringInteractionsTrust2021}.

These studies provide a theoretical context for
investigating emotional-attachment and social-avoidance
markers. However, synthetic pipeline results cannot
confirm or contradict empirical findings about
human participants.

\subsection{Limitations and Threats to Validity}
\label{subsec:limitations}

\textbf{Internal validity.}
The synthetic dataset lacks author identifiers.
Treating each observation as an independent author
is therefore a pipeline-testing assumption rather
than a verified sampling property. Human coding
reliability has not yet been established.

\textbf{External validity.}
The planned authentic study focuses on English-language
accounts concerning Replika and Character.AI.
Its future findings may not generalize to other
platforms, languages, or populations.

\textbf{Construct validity.}
Normalized textual marker frequencies are indirect
indicators of emotional attachment and social
avoidance. They cannot establish psychological
diagnoses or fully represent these constructs.

\textbf{Causal limitations.}
The proposed study is observational and
cross-sectional. Statistical associations would
not establish causal effects of AI-companion use.

\textbf{Data status.}
The current results validate computational
procedures only. Empirical conclusions require
authentic data, documented ethical safeguards,
verified sampling units, and independent coding
reliability assessment.
